% GENERATED FILE. Edit canonical lessons, then run npm run generate:books.
% canonical-source-hash: fnv1a64:b4fa0692313c6292
% canonical-lessons: PA-C63-labbhna, PA-C63-suttna, PA-C63-mit, PA-C63-machchhi, PA-C63-chaul

\chapter{To find, To throw, Meat, Fish, Rice}
\label{ch:pa-to-find-to-throw-meat-fish-rice}
\hlchaptermodality{63}

\begin{chapteropening}
\textbf{By the end of this chapter:} \emph{I can say to find, to look for, to throw, meat, a fish and rice.}

Complete the last lesson of chapter 63: i can say to find, to look for, to throw, meat, a fish and rice.
\end{chapteropening}

\section[labbhṇā]{\pa{ਲੱਭਣਾ} (labbhṇā) — to find, to look for}

\label{lesson:PA-C63-labbhna}



\begin{quote}
Before the new one: say the Punjabi for to leave, then the Punjabi for to wait for.
\end{quote}

\subsection*{You'll want to know: \pa{ਲੱਭਣਾ}}
\textbf{\pa{ਲੱਭਣਾ}} — \emph{labbhṇā} — \textquotedblleft{}to find, to look for\textquotedblright{}.

\begin{grammarlens}[title={What you've built}]
One more word for this chapter.
\end{grammarlens}

\subsection*{Your turn}
\begin{itemize}
\raggedright
  \item \emph{Say it:} \emph{labbhṇā}
  \item \emph{Say it:} \emph{labbhṇā}, once more
  \item \emph{Say it:} say \emph{labbhṇā}
  \item \emph{Recall:} say the Punjabi for to leave, then the Punjabi for to wait for, then say \emph{labbhṇā} again
\end{itemize}

\begin{tcolorbox}[breakable,colback=teal!4,colframe=teal!35!black,title={Before you move on}]
What does \pa{ਲੱਭਣਾ} mean? (\textquotedblleft{}To find, to look for\textquotedblright{}.) Say it once more.
\end{tcolorbox}
\section[suṭṭṇā]{\pa{ਸੁੱਟਣਾ} (suṭṭṇā) — to throw}

\label{lesson:PA-C63-suttna}



\begin{quote}
Before the new one: say the Punjabi for to wait for, then the Punjabi for to find.
\end{quote}

\subsection*{You'll want to know: \pa{ਸੁੱਟਣਾ}}
\textbf{\pa{ਸੁੱਟਣਾ}} — \emph{suṭṭṇā} — \textquotedblleft{}to throw\textquotedblright{}.

\begin{grammarlens}[title={What you've built}]
One more word for this chapter.
\end{grammarlens}

\subsection*{Your turn}
\begin{itemize}
\raggedright
  \item \emph{Say it:} \emph{suṭṭṇā}
  \item \emph{Say it:} \emph{suṭṭṇā}, once more
  \item \emph{Say it:} say \emph{suṭṭṇā}
  \item \emph{Recall:} say the Punjabi for to wait for, then the Punjabi for to find, then say \emph{suṭṭṇā} again
\end{itemize}

\begin{tcolorbox}[breakable,colback=teal!4,colframe=teal!35!black,title={Before you move on}]
What does \pa{ਸੁੱਟਣਾ} mean? (\textquotedblleft{}To throw\textquotedblright{}.) Say it once more.
\end{tcolorbox}
\section[mīṭ]{\pa{ਮੀਟ} (mīṭ) — meat}

\label{lesson:PA-C63-mit}



\begin{quote}
Before the new one: say the Punjabi for to find, then the Punjabi for to throw.
\end{quote}

\subsection*{You'll want to know: \pa{ਮੀਟ}}
\textbf{\pa{ਮੀਟ}} — \emph{mīṭ} — \textquotedblleft{}meat\textquotedblright{}.

\begin{grammarlens}[title={What you've built}]
One more word for this chapter.
\end{grammarlens}

\subsection*{Your turn}
\begin{itemize}
\raggedright
  \item \emph{Say it:} \emph{mīṭ}
  \item \emph{Say it:} \emph{mīṭ}, once more
  \item \emph{Say it:} say \emph{mīṭ}
  \item \emph{Recall:} say the Punjabi for to find, then the Punjabi for to throw, then say \emph{mīṭ} again
\end{itemize}

\begin{tcolorbox}[breakable,colback=teal!4,colframe=teal!35!black,title={Before you move on}]
What does \pa{ਮੀਟ} mean? (\textquotedblleft{}Meat\textquotedblright{}.) Say it once more.
\end{tcolorbox}
\section[machchhī]{\pa{ਮੱਛੀ} (machchhī) — a fish}

\label{lesson:PA-C63-machchhi}



\begin{quote}
Before the new one: say the Punjabi for to throw, then the Punjabi for meat.
\end{quote}

\subsection*{You'll want to know: \pa{ਮੱਛੀ}}
\textbf{\pa{ਮੱਛੀ}} — \emph{machchhī} — \textquotedblleft{}a fish\textquotedblright{}.

\begin{grammarlens}[title={What you've built}]
One more word for this chapter.
\end{grammarlens}

\subsection*{Your turn}
\begin{itemize}
\raggedright
  \item \emph{Say it:} \emph{machchhī}
  \item \emph{Say it:} \emph{machchhī}, once more
  \item \emph{Say it:} say \emph{machchhī}
  \item \emph{Recall:} say the Punjabi for to throw, then the Punjabi for meat, then say \emph{machchhī} again
\end{itemize}

\begin{tcolorbox}[breakable,colback=teal!4,colframe=teal!35!black,title={Before you move on}]
What does \pa{ਮੱਛੀ} mean? (\textquotedblleft{}A fish\textquotedblright{}.) Say it once more.
\end{tcolorbox}
\section[chaul]{\pa{ਚੌਲ} (chaul) — rice}

\label{lesson:PA-C63-chaul}



\begin{quote}
Before the new one: say the Punjabi for meat, then the Punjabi for a fish.
\end{quote}

\subsection*{You'll want to know: \pa{ਚੌਲ}}
\textbf{\pa{ਚੌਲ}} — \emph{chaul} — \textquotedblleft{}rice\textquotedblright{}.

\begin{grammarlens}[title={What you've built}]
One more word for this chapter.
\end{grammarlens}

\subsection*{Your turn}
\begin{itemize}
\raggedright
  \item \emph{Say it:} \emph{chaul}
  \item \emph{Say it:} \emph{chaul}, once more
  \item \emph{Say it:} say \emph{chaul}
  \item \emph{Recall:} say the Punjabi for meat, then the Punjabi for a fish, then say \emph{chaul} again
\end{itemize}

\begin{tcolorbox}[breakable,colback=teal!4,colframe=teal!35!black,title={Before you move on}]
What does \pa{ਚੌਲ} mean? (\textquotedblleft{}Rice\textquotedblright{}.) Say it once more.
\end{tcolorbox}
